% \begin{abstract}
% The recent cohort of publicly available generative models can produce human expert level content across a variety of topics and domains.
% Given a model in this cohort as a base model, methods such as parameter efficient fine-tuning, in-context learning, and constrained decoding have further increased generative capabilities and improved both computational and data efficiency.
% Entire collections of derivative models have emerged as a byproduct of these methods and each of these models has a set of associated covariates such as a score on a benchmark, an indicator for if the model has (or had) access to sensitive information, etc. that may or may not be available to the user.
% For some model-level covariates, it is possible to use “similar” models to predict an unknown covariate.
% In this paper we extend recent results related to embedding-based representations of generative models -- the \textit{data kernel perspective space} -- to classical statistical inference settings. 
% We demonstrate that using the perspective space as the basis of a notion of ``similar" is effective for multiple model-level inference tasks.
% \end{abstract}

\begin{abstract}
Generative models are capable of producing human-expert level content across a variety of topics and domains.
As the impact of generative models grows, it is necessary to develop statistical methods to understand collections of available models.
These methods are particularly important in settings where the user may not have access to information related to a model's pre-training data, weights, or other relevant model-level covariates.
In this paper we extend recent results on representations of black-box generative models to model-level statistical inference tasks.  
We demonstrate that the model-level representations are effective for multiple inference tasks.
\end{abstract}